\chapter{Espaces métriques}

\lecture{7 septembre}
\section{Rappel analyse avancée et algèbre linéaire}
\subsection{Métrique sur $\Rn$}
\begin{definition}[Distance euclidienne] \label{def: Distance euclidienne}
	On définit la \textbf{distance euclidienne} sur $\Rn$ comme
	\[
	\begin{array}{cccc}
		d_{euc} : & \Rn \times \Rn & \longrightarrow & \R \\
					& (v, w) & \longmapsto & \left( \sum_{i=1}^n (v_i - w_i)^2 \right)^\frac12
	\end{array}
	\]
	ayant les propriétés importantes suivantes :
	\begin{enumerate}
		\item $\forall u,v,w \in \Rn$, $d_{euc}(u,w) \leqslant d_{euc}(u,v) + d_{euc}(v,w)$
		
		\item $d_{eud}(u,v) = 0 \iff v = u$
		
		\item $\forall u,v \in \Rn, \quad d_{euc}(u,v) = d_{euc}(v,u)$
	\end{enumerate}
\end{definition}

\begin{definition}[Produit scalaire et norme euclidienne] \label{def: Produit scalaire et norme euclidienne}
	On définit le \textbf{produit scalaire euclidien} comme
	\[
	\begin{array}{cccc}
		\langle \cdot, \cdot \rangle_{euc} : & \Rn \times \Rn& \longrightarrow & \R \\
												& (u,v) & \longmapsto & \sum_{i=1}^n u_i v_i 
	\end{array}
	\]
	dont la \textbf{norme euclidienne} induite est
	\[
	\begin{array}{cccc}
		\norme{\cdot}_{euc} : & \Rn & \longrightarrow & \Rn \\
								& v & \longmapsto & \langle v,v \rangle^\frac12 = \left(\sum_{i=1}^n v_i^2 \right)^\frac12
	\end{array}
	\]
	On remarque alors que $d_{euv}(u,v) = \norme{u -v}$.
\end{definition}

\subsection{Continuité}
\begin{definition}[Continuité] \label{def: Continuité}
	Soit $f : \Rn \to \Rm$ une application. Soit $v \in \Rn$, on dit que $f$ est \textbf{continue} en $v$ si et seulement si
	\[
	\forall \varepsilon >0, ~\exists \delta > 0, ~\forall w \in \Rn, \quad d_{euc}(v,w) < \delta \implies d_{euc}(f(v), f(w)) <  \varepsilon
	\]
	Par extension, on dira que $f$ est continue si elle est continue en tout $v \in \Rn$.
\end{definition}

\begin{definition}[Boule ouverte euclidienne] \label{def: Boule ouverte euclidienne}
	Soit $v \in \Rn$ et $r > 0$, on définit la \textbf{boule ouverte euclidienne} autour de $v$, de rayon $r$ comme,
	\[
	B_{euc}(v,r) = \{w \in \Rn \mid d_{euc}(v,w) < r \} \subset \Rn
	\]
\end{definition}

\begin{theorem}[Caractérisation de la continuité] \label{thm: Caractérisation de la continuité}
	Soit $f : \Rn \to \Rm$ une application, alors $f$ est continue si et seulement si,
	\[
	\forall w \in \Rn, ~\forall r > 0, ~\forall v \in f^{-1}(B_{euc}(w,r)), ~\exists s > 0, \quad B_{euc}(v, s) \subset f^{-1}(B_{euc}(w,r))
	\]
	Autrement dit, $f$ est continue si et seulement si toute pré-image d'un ouvert par $f$ est un ouvert.
\end{theorem}

\prf{En premier lieu on remarque que pour $v \in \Rn$
	\[
	d_{euc}(f(v), f(w)) < \varepsilon \iff f(w) \in B_{euc}(f(v), \varepsilon) \et d_{euc}(v,w) < \delta \iff w \in B_{euc}(v, \delta)
	\]
	Ainsi on veut trouver $\delta > 0$ tel que $w \in B_{euc}(v, \delta) \implies f(w) \in B_{euc}(f(v), \varepsilon)$. Or
	\[
	f(w) \in B_{euc}(f(v), \varepsilon) \iff w \in f^{-1}(B_{euc}(f(v), \varepsilon))
	\]
	Or, par hypothèse, $\forall u \in f^{-1}(B_{euc}(f(v), \varepsilon))$ il existe $s > 0$ tel que
	\[
	B_{euc}(u,s) \subset f^{-1}(B_{euc}(f(v), \varepsilon))
	\]
	Et donc il suffit de prendre $\delta = s$, ce qui montre bien la continuité de $f$ en $v$.
	
	Réciproquement, soient $w \in \Rm, r > 0$ et $v \in f^{-1}(B_{euc}(w,r))$. Donc par définition de la pré-image, $f(v) \in B_{euc}(w,r)$, d'où $d_{euc}(w, f(v)) < r$. On pose $\varepsilon = r - d_{euc}(w, f(v)) > 0$, donc par continuité de $f$, il existe $\delta > 0$ tel que
	\[
	u \in B_{euc}(v, \delta) \implies f(u) \in B_{euc} (f(v),\varepsilon)
	\]	
	or
	\[
	d_{euc}(w, f(u)) \leqslant d_{euc}(w, f(v)) + d_{euc}(f(v), f(u)) < d_{euc}(w,f(v)) + \varepsilon = r
	\]
	donc $u \in f^{-1}(B_{euc}(w,r))$. Et donc $B_{euc}(v, \delta) \subset f^{-1}(B_{euc}(w,r))$ ainsi la première implication est bien vérifiée pour $s = \delta$.
}

\section{Généralisation de la métrique}
\begin{definition}[Espace métrique] \label{def: Espace métrique}
	Soit $X$ un ensemble, une \textbf{métrique} sur $X$ est une application $d : X \times X \to \R$ telle que
	\begin{enumerate}
		\item $d(x,x') = 0 \iff x = x'$
		\item $\forall x,x' \in X, \quad d(x,x') = d(x', x)$
		\item $\forall x,x',x'' \in X, \quad d(x,x'') \leqslant d(x,x') + d(x',x'')$
	\end{enumerate}
	Ainsi on définit un \textbf{espace métrique} comme le couple $(X,d)$. De plus si $X$ est fini, on parlera d'espace métrique \textbf{fini}. Et fininalement s'il existe $M \in \R$ tel que 
	\[
	\forall x, x' \in \R, \quad d(x,x') \leqslant M
	\]
	alors on dira que $(X,d)$ est un espace métrique \textbf{borné}.
\end{definition}

\begin{lemme}[Positivité] \label{Positivité}
	Soit $(X,d)$ un espace métrique. Alors 
	\[
	\forall x,x' \in X, d(x,x') \geqslant 0
	\]
\end{lemme}
\prf{On a simplement,
\[
\forall x, x' \in X, \quad 0 = d(x,x) \leqslant d(x,x') + d(x,x') = 2d(x,x') \implies d(x,x') \geqslant 0
\]
}

\begin{exemple}[Espace métrique usuel]
		Soit $V$ un $\R$-espace vectoriel muni d'un produit scalaire $\langle \cdot,\cdot \rangle$, alors 
		\[
		\begin{array}{cccc}
		d_{\langle \cdot, \cdot \rangle} :& V \times V &\longrightarrow &\R \\ &(v,w) &\longmapsto &\norme{v - w} = \langle v -w, v-w \rangle^\frac12
		\end{array}
		\]
		est une métrique sur $V$.

\end{exemple}

\begin{exemple}[Métrique de chemin]
	Soit $G= (E,V)$ un graphe simple, alors on peut définir la \textbf{métrique de chemin} sur $V$ déterminée $G$ pour tout $v,w \in V$ par
	\[
	d_G(v,w) = \min \{ n \mid \text{il existe un chemin $(v_0, \ldots, v_n)$ de $v$ à $w$}\}
	\]
	où un \textbf{chemin} de $v$ à $w$ consiste en une suite $(v_0, \ldots, v_n)$ de sommets telle que $v_0 = v$ et $v_n = w$, et $\{v_i, v_{i+1}\} \in E$ pour tout $i \in \llbracket 0, n-1 \rrbracket$.	
\end{exemple}

\lecture{14 septembre}


\begin{exemple}[Métrique de Hamming] \label{ex: Métrique de Hamming}
Soit $A$ un ensemble, et soit $n \geqslant 0$ un nombre naturel, alors on pose
\[
W_n(A) = \{ (a_1, \ldots, a_n) \mid a_i \in A \}
\]	
l'ensemble des mots de longueur $n$ dans l'alphabet $A$. Alors la \emph{métrique de Hamming} sur $W_n(A)$ est définie par
\[
\begin{array}{cccc}
	d_H : & W_n(A) \times W_n(A) & \longrightarrow & \R \\
			& (a,b) & \longmapsto & \# \{ i \mid a_i \neq b_i \}
\end{array}
\]
On vérifie facilement que c'est bien une métrique, de plus on remarque que $(W_n(A), d_h)$ est un espace métrique borné.
\end{exemple}

\begin{exemple}[Métrique uniforme]
	Soit $(X,d)$ un espace métrique et soit $B$ un ensemble quelconque, alors on pose 
	\[
	\mathcal F(B, X) = \{ f : B \to X\}
	\]
	l'ensemble des applications de $B$ à valeurs dans $X$. Si $(X,d)$ est borné, alors la \textbf{métrique uniforme} définie par
	\[
	d_{unif}(f,g) = \sup_{b \in B} d(f(b), g(b))
	\]
	existe.
\end{exemple}

\section{Construction d'espaces métriques}
On se demande comment construire un espace métrique à partir d'autres espaces métriques donnés.

\subsection{Sous-espace métrique}

\begin{definition}[Sous-espace métrique] \label{def: Sous-espace métrique}
	Soit $(X,d)$ un espace métrique et $Y \subset X$ un sous-ensemble de $X$. Si on pose
	\[
	 \begin{array}{cccc}
	 	d_Y = d_{\mid Y \times Y} : & Y \times Y & \longrightarrow & \R \\
	 								& (y_1, y_2) & \longmapsto & d(y_1, y_2)
	 \end{array}
	\]
	alors $(Y, d_Y)$ est appelé le \textbf{sous-espace métrique} déterminé par $Y$.
\end{definition}

\begin{exemple}
	En reprenant l'\autoref{ex: Métrique de Hamming} alors si on pose
	\[
	\mathrm{Inj}_n(A) = \{ (a_1, \ldots a_n) \in W_n(A) \mid a_i \neq a_j \forall i\neq j\} \subset W_n(A)
	\]
	l'ensemble des mots injectifs, alors $(\mathrm{Inj}_n(A), d_h)$ est un sous-espace métrique.
\end{exemple}

\subsection{Métrique produit}
Si on souhaire définir une métrique sur $X_1 \times \ldots \times X_n$, au lieu d'avoir
\[
d_\pi : (X_1 \times \ldots \times X_n) \times (X_1 \times \ldots \times X_n) \longrightarrow \R
\]
comme une application partant de deux grands produit, on peut le considérer comme $n$ produit de $X_i$ avec lui-même.
\begin{figure}[H]
	\centering
	\begin{tikzcd}
		{(X_1 \times \ldots \times X_n) \times (X_1 \times \ldots \times X_n)} && {\mathbb R} \\
		{(X_1 \times X_1) \times \ldots \times (X_n \times X_n)} \\
		{\mathbb R^n}
		\arrow[from=1-1, to=1-3]
		\arrow["\simeq", from=1-1, to=2-1]
		\arrow["{d_1 \times \ldots \times d_n}"', from=2-1, to=3-1]
		\arrow["{\| \cdot \|_{euc}}"{description}, bend right=30, from=3-1, to=1-3]
	\end{tikzcd}
	\caption{Diagramme commutatif}
\end{figure}
Et donc on a naturellement la définition suivante pour définir une métrique produit induite par les $X_i$.
\begin{definition}
	Soient $(X_1, d_1), (X_2, d_2), \ldots (X_n, d_n)$ des espaces métriques.  Alors  la \textbf{métrique produit} induite sur $X_1 \times X_2 \times \ldots X_n$ est
	\[
	\begin{array}{cccc}
		d_\pi : & (X_1, \times \ldots \times X_n) \times (X_1 \times \ldots \times X_n) & \longrightarrow & \R \\
				& ((x_1, \ldots , x_n), (x_1', \ldots, x_n')) & \longmapsto & \left(\sum_{i=1}^n d_i(x_i, x_i')^2\right)^\frac12
	\end{array} 
	\]
	Il est ensuite facile de vérifier que $d_\pi$ définit bien une métrique.
\end{definition}

\section{Application entre espaces métriques}
\begin{definition}[Fonction préservant les distances et isométries]
	Soit $(X, d_X)$ et $(Y, d_Y)$ deux espaces métriques et soit $f : X \to Y$ une application. Alors on dit que $f$ \textbf{préserve les distances} si
	\[
	\forall x_1, x_2 \in X, \quad d_X(x_1, x_2) = d_Y(f(x_1), f(x_2))
	\]
	De plus on dira que $f$ est une \textbf{isométrie} si elle préserve les distances et est surjective (et donc bijective).
\end{definition}

\begin{remarque}
	Il est facile de voir que si $f$ préserve les distances alors elle est nécessairement injective, en effet on a
	\[
	f(x_1) = f(x_2) \implies 0 = d_Y(f(x_1), f(x_2)) = d_X(x_1, x_2) \implies x_1 = x_2
	\]
	d'où l'injectivité.
\end{remarque}

\begin{definition}
	Soient $X, d_X)$ et $(Y, d_Y)$ des espaces métriques et $f : X \to Y$ une application. Soit encore $K \in \R_{> 0}$. On dit que $f$ est $K$\textbf{-lipschitzienne} par rapport à $d_X$ et $d_Y$ si
	\[
	\forall x_1, x_2 \in X, \quad d_Y(f(x_1), f(x_2)) \leqslant K d_X(x_1, x_2)
	\]

\end{definition}
Il reste une dernière notion qui utilise la notion de la continuité. Ainsi on a
\begin{definition}
	Soient $X, d_X)$ et $(Y, d_Y)$ des espaces métriques et $f : X \to Y$ une application. Alors on dit que $f$ est continue par rapport $d_X$ et $d_Y$ si 
	\[
	\forall \varepsilon > 0, \, \forall x \in X, \, \exists \delta >  0, \, \forall y \in X, \quad d_X(x,y) < \delta \implies d_Y(f(x), f(y)) < \varepsilon
	\]
\end{definition}	

\begin{remarque}
	Il est facile de voir que 
	\[
	\text{$f$ préservant les distances} \implies \text{$f$ Lipschitzienne} \implies \text{$f$ continue}
	\]
	Mais les réciproques ne sont généralement pas vraie, et donc la notion de préservation des distances et une notion beaucoup plus contraignante sur $f$, mais c'est en contre parti la notion la plus forte. De plus pour la continuité, on a toujours la caractérisation qui est analogue à la \hyperref[thm: Caractérisation de la continuité]{caractérisation de la continuité dans $\Rn$}.
\end{remarque}

\begin{theorem}
	Une application
	\[
	f : (X, d_X) \longrightarrow (Y, d_Y)
	\]
	est \textbf{continue} si et seulement si pour tout $y \in Y$, $r > 0$ et $x \in f^{-1}(B_{d_Y}(y,r))$, il existe un $\delta > 0$ tel que $B_{d_X} (x, \delta) \subset f^{-1} (B_{d_Y}(y,r))$.
\end{theorem}

\prf{La preuve est analogue au cas euclidien.}
